\section{What improves or worsens the plasma: evaluation of the dataset}\label{sec:qi_einfluss}

\textbf{Reading aid.} The QI quality (quasi-isodynamic) measures how well a plasma confines particles in the magnetic
field; a smaller value is better. The \emph{rank correlation} only says whether two quantities rise together (positive)
or in opposite directions (negative). It is \emph{not a cause}: we never changed a plasma and measured the effect.

The user asked what has a positive and what a negative effect on the plasma. The plasma shape comes from the public
Proxima dataset; our chain only evaluates coils and support structure for it.

\subsection*{Expectation beforehand}
Among the 182 pre-selected coil sets (force scan) the aspect ratio was the strongest influence (+0.30, larger is worse).
The same sign was expected in the full dataset, with a larger magnitude (0.3 to 0.5).

\subsection*{Result}
158,685 of 182,222 plasmas have a valid QI. Rank correlation with $\log_{10}$(QI):

\begin{tabular}{lr}
\hline
Quantity & Rank correlation \\ \hline
edge rotational transform (per period) & $+0.35$ \\
minimum normalised gradient length & $-0.34$ \\
aspect ratio & $-0.31$ \\
maximum elongation & $-0.30$ \\
vacuum well & $-0.30$ \\
edge mirror ratio & $+0.27$ \\
average triangularity & $+0.22$ \\
axis mirror ratio & $-0.21$ \\
flux compression in bad curvature & $-0.20$ \\
field periods & $+0.14$ \\ \hline
\end{tabular}

\smallskip
\noindent Median QI by field periods: 1: 0.0109; 2: 0.0085; 3: 0.0204; 4: 0.0145; 5: 0.0179. The best single plasma has
2 periods (QI $6\cdot10^{-5}$).

\subsection*{Comparison with the expectation: missed}
The sign for the aspect ratio \emph{reverses} ($-0.31$ instead of $+0.30$). Checking within aspect-ratio ranges shows
why: below 10 a larger aspect ratio is better (about $-0.17$ to $-0.34$), above 10 worse ($+0.25$ and more). Our 182 sets
all lie between 9 and 11, where the correlation has the other sign. The field periods are misleading as well:
the overall value ($+0.14$) hides $-0.14$ to $-0.37$ inside each period count.

\subsection*{Interpretation and limits}
\begin{itemize}\itemsep1pt
\item The correlations are weak (at most 0.35) and depend on the range considered. They give \emph{no} rule such as
  ``flatter is better''.
\item Only this holds: the best plasmas have few field periods (2); 3 periods are worst in the median.
\item The earlier statement that more force goes with better QI ($-0.26$) applies to the 182 sets only and has not been
  checked on the full dataset.
\item QI measures the field shape. Confinement and stability of the plasma are not calculated.
\end{itemize}
Source: \texttt{stellarator/qi\_einfluss.py}, result \texttt{docs/recherche/qi\_tragwerk/qi\_einfluss.md}.

\subsection*{Addendum: the quantities influence each other}
The user suspected that these are ranges that influence each other. A joint model (gradient boosting, checked on held-out
data) confirms this: it explains the QI quality with $R^2=0.90$, every single correlation at most $0.12$. The two mirror
ratios (on the axis and at the edge) matter most; individually they had opposite signs and thereby cancelled.
\emph{Reading aid:} the mirror ratio says how strongly the field varies along a field line; for good QI the
variations are similar on all surfaces.

Ranges of the best 1\,\% (QI at most $0.00084$; 5th to 95th percentile, all plasmas in brackets):
aspect ratio 7.5 to 9.7 (5.3 to 10.6); elongation 5.3 to 7.4; axis mirror ratio 0.28 to 0.40 (0.12 to 0.58),
edge 0.31 to 0.43 (0.20 to 0.85); edge rotational transform 0.07 to 0.24; at most 4 field periods.

\emph{Expectation beforehand:} $R^2>0.5$ (met), most important quantity rotational transform or gradient length
(missed). \emph{Limit:} the mirror ratios are themselves measures of the field shape and closely related to QI, so
they are no freely selectable design levers; it remains a correlation. Source: \texttt{stellarator/qi\_bereiche.py}.
